% !TEX root = ../../proposal.tex

\section{Compositionality in Diffusion Models}
\label{sec:rw_compositional}

The idea of composing diffusion models at inference time without retraining originated
with Liu et al.~\citep{liu2022compositional}, who showed that adding per-concept
classifier-free guidance terms gives an AND operator over concept conditions and a NOT
operator for negation. Concurrent work by Nie et al.~\citep{nie2021controllable} and
Du et al.~\citep{NEURIPS2020_49856ed4} approached the same problem from an
energy-based perspective, treating the composed distribution as a product of energy
functions. Du et al.~\citep{du2024compositional} later demonstrated that a
\emph{single} pretrained model suffices for multi-concept composition, replacing the
need for independently trained experts. Their analysis shows that dynamic guidance
schedules and adaptive step sizes mitigate the mode collapse and instability that
arise when concept scores are na\"{i}vely added, directly motivating the trajectory
dynamics analysis in Phase~1. Reduce, Reuse, Recycle~\citep{reduce_reuse_recycle_icml2023}
extended compositional diffusion to few-shot concept learning, showing that individual
concept encodings can be recombined without joint training data. These works collectively
establish score-level composition as a practical paradigm, but they do not theorise when
it should fail or characterise the failure modes empirically across structured concept
regimes. This gap is what Phase~1 addresses.

The theoretical account of why composition fails was provided by
Bradley~\citep{mechanisms2025projective}, who derived the Factorised Conditionals
condition: composition is well-defined if and only if the conditional score fields of the
composed concepts act on approximately orthogonal directions in feature space. When this
condition is violated --- either because concepts occupy the same feature dimensions or
because their joint configuration is absent from the training distribution --- the
composed score is incoherent and sampling diverges from the data manifold. Bradley's
framework predicts composition success from the geometry of concept representations,
not from the choice of composition operator. Our work is most directly related to this
theoretical programme: Phase~1 takes Bradley's conditions as a taxonomy scaffold,
operationalises them into four empirically testable concept groups, and measures for the
first time how the composability gap varies systematically across those groups. Bradley
identifies the structural conditions but does not study how failure modes differ across
regime types in practice, nor does the paper test whether better representations close
the gap. Phases~2 and~3 address both of these open questions.

Skreta et al.~\citep{skreta2024superposition} proposed Superposition Diffusion, which
combines multiple diffusion processes through a shared latent superposition rather than
score addition, targeting molecular and protein design. While this extends composition
beyond image generation, it does not address the representational account of compositional
failure that motivates our work. Our contribution focuses on \emph{when} and \emph{why}
inference-time composition fails and whether learned latent factorisation can remedy it.
